% !TeX program = xelatex
\documentclass[12pt, a4paper, titlepage]{article}
\usepackage{xeCJK}
\usepackage{amsmath, amssymb}
\usepackage{geometry}
\usepackage{longtable}
\usepackage{booktabs}
\usepackage{tabularx}
\usepackage{caption}
\usepackage{silence}
\WarningFilter{caption}{Unknown document class}
\WarningFilter{xeCJK}{Redefining CJKfamily}
\WarningFilter{relsize}{Failed to get list}

% -- Fonts --
\setmainfont{Times New Roman}
\setCJKmainfont{Hiragino Mincho ProN}

\geometry{reset, margin=2.5cm}

\title{LS vs CLS 保存性比較数値実験 仕様書（動的非一様格子・DCCD対応版）}
\author{}
\date{}

\begin{document}

\maketitle

\section*{本実験の目的}
本実験の目的は、急峻な勾配を持つ界面において、標準的な Level Set (LS) 法と Conservative Level Set (CLS) 法の保存性を比較することである。特に、界面の位置に合わせて格子を逐次リフレッシュする動的非一様格子環境下で、高精度スキーム (DCCD) とメトリクス評価が保存性に与える影響を定量化する。

\section*{1. 数理モデルと離散化スキーム}

\subsection*{1.1 支配方程式（1次元移流方程式と格子速度）}
格子が時間的に変化するため、計算面 $\xi$ における移流方程式は格子速度 $u_g = \frac{\partial x}{\partial t}$ を考慮した対偶形式となる：
$$\frac{\partial \psi}{\partial t} + (u - u_g) \frac{\partial \xi}{\partial x} \frac{\partial \psi}{\partial \xi} = 0$$ここで、$\frac{\partial \xi}{\partial x}$ は CCD 法を用いて算出される時空間メトリクスである。

\subsection*{1.2 滑らかなヘヴィサイド関数と界面幅}$$\psi(x, t) = \frac{1}{2} \left[ 1 + \tanh\left( \frac{\phi(x, t)}{2\epsilon} \right) \right]$$動的格子においては、局所的な最小格子間隔 $h_{\min}$ に基づいて $\epsilon$ を一定に保つか、あるいは格子解像度に合わせて $\epsilon$ をスケールさせる。

\subsection*{1.3 空間離散化}
\begin{itemize}
    \item DCCD (Dissipative Combined Compact Difference): 移流項に対して 6 次精度を維持しつつ、格子リフレッシュに伴う微細な数値振動を抑制する。
    \item CCD メトリクス評価: 非一様格子の幾何学的性質（$\frac{\partial x}{\partial \xi}$ 等）を 6 次精度で評価し、幾何学的保存則（GCL）を満足させる。
\end{itemize}

\section*{2. 動的格子リフレッシュ仕様}

\subsection*{2.1 格子生成関数}
界面位置 $x_c(t)$ を中心とした密度関数 $w(x)$ を定義し、代数的格子生成または微分方程式法により格子 $x_i(t)$ を更新する。
\begin{itemize}
    \item リフレッシュ頻度: 毎タイムステップ、または界面が細密領域を脱する前。
    \item 格子集中化: 界面近傍の格子解像度を背景格子の 4〜8 倍程度に設定。
\end{itemize}

\subsection*{2.2 物理量の再マッピング (Remapping)}
格子更新後、旧格子上の値 $\psi_{\text{old}}$ を新格子 $\psi_{\text{new}}$ へ転送する：
\begin{itemize}
    \item LS法: 3次以上のスプラインまたは CCD 補間。
    \item CLS法: 積分型保存特性を維持するため、コントロールボリューム間のフラックスを考慮した保存補間、または再初期化ステップを介した補正を行う。
\end{itemize}

\section*{3. 数値計算条件}

\noindent\begin{tabularx}{\linewidth}{|l|X|X|}
\hline
項目 & 設定値 & 備考 \\
\hline
格子構造 & 動的非一様格子 & 界面位置 $x_c$ に追従してリフレッシュ \\
格子数 $N$ & 128（実質解像度は界面付近で 512 相当） &  \\
移流スキーム & DCCD (6次) & 散逸型コンパクト差分 \\
再初期化 & CLS 圧縮・拡散ステップ & 格子更新後のプロファイル整形 \\
時間積分 & TVD-RK3 &  \\
\hline
\end{tabularx}

\section*{4. 実験手順}
\begin{enumerate}
    \item 初期格子生成: 界面 $(x=0.25)$ 付近を細密化した格子を作成。
    \item 移流と格子追従: 界面の移動量に応じて格子系全体をシフト、または再構成する。
    \item メトリクス更新: 新しい格子配置における $\frac{\partial \xi}{\partial x}$ を CCD で再計算。
    \item 保存性の計測: 領域全体を 1 周移流させた後の総質量の変化を記録。
\end{enumerate}

\section*{5. 評価指標}

\subsection*{5.1 質量保存誤差 ($E_{\text{mass}}$)}$$E_{\text{mass}}(t) = \frac{|\int \psi(x, t) \sqrt{g} d\xi - M_0|}{M_0}$$
$\sqrt{g}$ は格子ヤコビアン。格子リフレッシュ時の補間誤差がここに蓄積される。

\subsection*{5.2 形状保持特性}
界面の勾配（厚み $\epsilon$）が、格子リフレッシュを繰り返す過程で拡散（増大）または急峻化（減少）していないかを評価。

\section*{6. 期待される結果}
\begin{itemize}
    \item 格子リフレッシュの影響: 単純な LS 法では、補間を繰り返すごとに質量の散逸が加速する。
    \item CLS + DCCD の優位性: 保存形形式と高精度なメトリクス評価により、格子が動的に変化し続けても全質量が極めて高い精度で維持される。
    \item DCCD の安定性: 急峻な変化部が格子間を移動する際、非物理的な振動（オーバーシュート）が DCCD によって効果的に抑制される。
\end{itemize}

\end{document}